\scanpage{222}
\mspage{L 19}
% archive sheet number: 221
Récoltes et Semailles j'ai pu me rendre compte, et de façon saisissante en ces tout derniers mois surtout (avec l'écriture des Quatre Opérations\handwritten{)}, qu'il y a eu depuis mon départ de la scène mathématique une stupéfiante \mbox{d\,é\,g\,r\,a\,d\,a\,t\,i\,o\,n} dans l'esprit qui aujourd'hui fait loi dans les milieux qqe j'avais connus, et (me semble-t-il, dans une large mesure au moins) dans le monde mathématique en général (\textsuperscript{\unsure{x|*}}). Il est possible même,\deleted{\unsure{que ma}}ntant par ma personnalité mathématique très particulière que par les conditions qui ont entouré mon départ, que celui-ci ait agi comme un catalyseur dans une évolution qui était déjà en train de se faire (\textsuperscript{\unsure{xx|**}}) - une évolution dont je n'ai \inserted{alors} rien su percevoir (pas plus qu'aucun autre de m\unsure{e|é}s collègues et amis, à la seule exception peut-être de Claude Chevalley\unsure{).} L'aspect de cette dégradation auquel je pense surtout ici (qui \inserted{\handwritten{en}} est \inserted{\handwritten{juste}} \uline{un} aspect parmi de nombreux autres (\textsuperscript{\unsure{xxx|***}})) est le \mbox{m\,é\,p\,r\,i\,s} \mbox{t\,a\,c\,i\,t\,e} , quand ce n'est la dérision sans équivoque, à l'encontre de ce qui (en mathématique, en l'occurence) ne s'apparente pas au pur travail du marteau sur l'enclume ou sur le burin - le mépris des processus créateurs les plus délicats (et \deleted{parf} souvent de moindre apparence\handwritten{)}, de tout ce qui est \deleted{inspiration,} \mbox{i\,n\,s\,p\,i\,r\,a\,t\,i\,o\,n} , \mbox{r\,ê\,v\,e} , \mbox{v\,i\,s\,i\,o\,n} (si puissantes et si fertiles soient-elles), et même (à la limite) de toute \mbox{i\,d\,é\,e} , si clairement conçue et formulée soit-elle : de tout ce qui n'est écrit et \mbox{p\,u\,b\,l\,i\,é} noir sur blanc, sous forme d'énon\deleted{\unsure{n}}cés purs et durs, répertoriables et répertoriés, mûrs pour les \char34{}banques de données\char34{} \deleted{\unsure{d'|s'}}engouffrés dans les inépuisables mémoires de nos mégaordinateurs.

Il y a eu (pour reprendre une expression de C.L. Siegel (\textsuperscript{\unsure{xxxx|****}}) un extraordinaire \deleted{\unsure{\char34{}aplatissement\char34{}, une}}\inserted{\char34{}\mbox{a\,p\,l\,a\,t\,\unsure{i|s}\,s\,s\,e\,m\,e\,n\,t}\char34{}}, un \char34{}\mbox{r\,é\,t\,r\,\unsure{é|e}\,c\,i\,s\,s\,e\,m\,e\,n\,t}\char34{} de la pensée mathématique,\deleted{\unsure{à|x}} dépouillée d'une dimension essentielle, de tout son \char34{}versant d'ombre\char34{} , du versant \char34{}féminin\char34{}. Il est vrai que par une tradition ancestrale, ce versant-là du travail de découverte restait dans une large mesure\handwritten{/}occulté, personne (autant dire) n'en \mbox{p\,a\,r\,l\,a\,i\,t} jamais - mais le contact \inserted{\handwritten{vivant}} avec les sources profondes du rêve, qui alimentent les grandes visions et les grands desseins, n'avait jamais encore (à ma connaissance) été perdu. Il semblerait qu\deleted{e}\handwritten{e} \inserted{\handwritten{dès à}} présent nous soyons déjà entrés dans une \mbox{é\,p\,o\,q\,u\,e} \mbox{d\,e} \mbox{d\,e\,s\,s\,è\,c\,h\,e}\deleted{\unsure{men}}\mbox{m\,e\,n\,t}, où cette source est, non point tarie certes, mais où l'accès à elle est condamné\deleted{\unsure{e}}, par le verdict sans appel du mépris général et par les représailles de la dérision. \handwritten{\(\big\Vert\)}Nous voilà approcher du moment, semble-t-il, où sera éradiqué en chacun non seulement le \mbox{s\,o\,u\,v\,e\,n\,i\,r} de tout travail proche de la source, du travail \char34{}au féminin\deleted{\unsure{x}}\char34{} (ridiculisé
